% Part: first-order-logic
% Chapter: syntax-and-semantics
% Section: terms-formulas

\documentclass[../../../include/open-logic-section]{subfiles}

\begin{document}

\olfileid{fol}{syn}{frm}

\olsection{Termen en \printtoken{P}{formula}}

Zodra een eerste-ordetaal~$\Lang L$ is gegeven, kunnen we uitdrukkingen
definiëren die uit het basisvocabulaire van~$\Lang L$ zijn opgebouwd.
Daartoe behoren in het bijzonder \emph{termen} en \emph{!!{formula}s}.

\begin{defn}[Termen]
\ollabel{defn:terms}
De verzameling \emph{termen}~$\Trm[L]$ van~$\Lang L$ wordt inductief
gedefinieerd door:
\begin{enumerate}
\item Iedere !!{variable} is een term.
\item Iedere !!{constant} van~$\Lang L$ is een term.
\item Als $f$ een $n$-plaatsige !!{function} is en $t_1$, \dots, $t_n$
  termen zijn, dan is $\Atom{f}{t_1, \ldots, t_n}$ een term.
\tagitem{limitClause}{
  Niets anders is een term.}{}
\end{enumerate}
Een term die geen !!{variable}s bevat, heet een \emph{gesloten term}.
\end{defn}

\begin{explain}
De !!{constant}s treden in onze specificatie van de taal en van de
termen op als een afzonderlijke categorie symbolen, maar hadden ook als
nulplaatsige !!{function}s kunnen worden opgenomen. De tweede clausule
van de termdefinitie zou dan overbodig zijn. We hoeven
$\Atom{f}{t_1, \ldots, t_n}$ slechts als $f$ zelf op te vatten als $n =
0$.
\end{explain}

\begin{defn}[Formules]
\ollabel{defn:formulas}
De verzameling \emph{!!{formula}s}~$\Frm[L]$ van de taal~$\Lang L$
wordt als volgt inductief gedefinieerd:
\begin{enumerate}
\tagitem{prvFalse}{$\lfalse$ is een atomaire !!{formula}.}{}

\tagitem{prvTrue}{$\ltrue$ is een atomaire !!{formula}.}{}

\item Als $R$ een $n$-plaatsig !!{predicate} van~$\Lang L$ is en $t_1$, \dots,
  $t_n$ termen van~$\Lang L$ zijn, dan is $\Atom{R}{t_1,\ldots, t_n}$ een
  atomaire !!{formula}.

\item Als $t_1$ en $t_2$ termen van~$\Lang L$ zijn, dan is $\Atom{\eq}{t_1, t_2}$
  een atomaire !!{formula}.

\tagitem{prvNot}{Als $!A$ !!a{formula} is, dan is $\lnot !A$
  !!a{formula}.}{}

\tagitem{prvAnd}{Als $!A$ en $!B$ !!{formula}s zijn, dan is $(!A \land
  !B)$ !!a{formula}.}{}

\tagitem{prvOr}{Als $!A$ en $!B$ !!{formula}s zijn, dan is $(!A \lor !B)$
  !!a{formula}.}{}

\tagitem{prvIf}{Als $!A$ en $!B$ !!{formula}s zijn, dan is $(!A \lif !B)$
  !!a{formula}.}{}

\tagitem{prvIff}{Als $!A$ en $!B$ !!{formula}s zijn, dan is $(!A \liff !B)$
  !!a{formula}.}{}

\tagitem{prvAll}{Als $!A$ !!a{formula} is en $x$ !!a{variable},
  dan is $\lforall[x][!A]$ !!a{formula}.}{}

\tagitem{prvEx}{Als $!A$ !!a{formula} is en $x$ !!a{variable},
  dan is $\lexists[x][!A]$ !!a{formula}.}{}

\tagitem{limitClause}{Niets anders is !!a{formula}.}{}
\end{enumerate}
\end{defn}

\begin{explain}
De definities van de verzameling termen en die van de !!{formula}s zijn
\emph{inductieve definities}. In wezen construeren we de verzameling
!!{formula}s in oneindig veel stadia. In het beginstadium verklaren we
alle atomaire formules tot formules; dit komt overeen met de eerste
gevallen van de definitie, namelijk de gevallen voor
\iftag{prvTrue}{$\ltrue$, }{}%
\iftag{prvFalse}{$\lfalse$, }{}%
$\Atom{R}{t_1,\dots,t_n}$ en $\Atom{\eq}{t_1,t_2}$. Met ``atomaire
!!{formula}'' wordt dus iedere !!{formula} van deze vorm bedoeld.

De overige gevallen van de definitie geven regels om uit reeds
geconstrueerde !!{formula}s nieuwe !!{formula}s te construeren. In het
tweede stadium kunnen we daarmee uit atomaire !!{formula}s nieuwe
!!{formula}s construeren. In het derde stadium construeren we nieuwe
formules uit de atomaire formules en de in het tweede stadium verkregen
formules, enzovoort. Een !!{formula} is precies wat uiteindelijk in een
dergelijk stadium wordt geconstrueerd.
\end{explain}

Volgens afspraak schrijven we $\eq$ tussen zijn argumenten en laten we
de haakjes weg: $\eq[t_1][t_2]$ is een afkorting van
$\Atom{\eq}{t_1,t_2}$. Verder wordt $\lnot \Atom{\eq}{t_1,t_2}$
afgekort als $\eq/[t_1][t_2]$. Bij een formule $(!B \ast !C)$ die uit
$!B$, $!C$ met een tweeplaatsig connectief~$\ast$ is geconstrueerd,
laten we het buitenste paar haakjes vaak weg en schrijven we eenvoudig
$!B \ast !C$.

\begin{intro}
Sommige logicaboeken eisen dat de !!{variable}~$x$ in~$!A$ voorkomt
voordat we
\iftag{prvEx}{$\lexists[x][!A]$ }{}%
\iftag{notprvEx,notprvAll}{}{en }%
\iftag{prvAll}{$\lforall[x][!A]$ }{}%
rekenen tot
\iftag{notprvEx,notprvAll}{!!a{formula}}{!!{formula}s}.
Zonder deze eis ontstaan geen problemen en wordt de behandeling
eenvoudiger.
\end{intro}

\begin{tagblock}{defNot,defOr,defAnd,defIf,defIff,defTrue,defFalse,defEx,defAll}
\begin{defn}
Formules die met de gedefinieerde operatoren zijn geconstrueerd, moeten
als volgt worden opgevat:

\begin{tagenumerate}{defTrue,defFalse,defNot,defOr,defAnd,defIf,defIff,defEx,defAll}

\tagitem{defTrue}{$\ltrue$ is een afkorting van
  \iftag{prvFalse}{$\lnot\lfalse$}{$(!A \lor \lnot !A)$ voor een
    zekere vaste atomaire !!{formula}~$!A$}.}{}

\tagitem{defFalse}{$\lfalse$ is een afkorting van
  \iftag{prvTrue}{$\lnot\ltrue$}{$(!A \land \lnot !A)$ voor een
    zekere vaste atomaire !!{formula}~$!A$}.}{}

\tagitem{defNot}{$\lnot !A$ is een afkorting van $!A \lif \lfalse$.}{}

\tagitem{defOr}{$!A \lor !B$ is een afkorting van
  \iftag{prvAnd}{$\lnot(\lnot !A \land \lnot !B)$}{$\lnot !A \lif
    !B$}.}{}

\tagitem{defAnd}{$!A \land !B$ is een afkorting van
  \iftag{prvOr}{$\lnot(\lnot !A \lor \lnot !B)$}{$\lnot (!A \lif
    \lnot !B)$}.}{}

\tagitem{defIf}{$!A \lif !B$ is een afkorting van
  \iftag{prvOr}{$\lnot !A \lor !B)$}{$\lnot (!A \land \lnot !B)$}.}{}

\tagitem{defIff}{$!A \liff !B$ is een afkorting van $(!A \lif !B) \land (!B
  \lif !A)$.}{}

\tagitem{defAll}{$\lforall[x][!A]$ is een afkorting van $\lnot\lexists[x][\lnot !A]$.}{}

\tagitem{defEx}{$\lexists[x][!A]$ is een afkorting van $\lnot\lforall[x][\lnot !A]$.}{}
\end{tagenumerate}
\end{defn}
\end{tagblock}

Wanneer we met een taal voor een bepaalde toepassing werken, schrijven
we tweeplaatsige !!{predicate}s en !!{function}s vaak tussen de
betreffende termen, bijvoorbeeld $t_1 < t_2$ en $(t_1 + t_2)$ in de
taal van de rekenkunde en $t_1 \in t_2$ in de taal van de
verzamelingenleer. De opvolgerfunctie van de rekenkundige taal wordt
volgens afspraak zelfs \emph{na} haar argument geschreven:~$t'$.
Formeel zijn dit echter slechts conventionele afkortingen van
$\Atom{\Obj A^2_0}{t_1, t_2}$, $\Obj f^2_0(t_1, t_2)$,
$\Atom{\Obj A^2_0}{t_1, t_2}$ respectievelijk $f^1_0(t)$.

\begin{defn}[Syntactische identiteit]
Het symbool $\ident$ drukt syntactische identiteit tussen tekenreeksen
uit, dat wil zeggen: $!A \ident !B$ dan en slechts dan als $!A$ en $!B$
tekenreeksen van dezelfde lengte zijn die op iedere positie hetzelfde
symbool bevatten.
\end{defn}

Het symbool $\ident$ mag aan weerszijden staan van door concatenatie
verkregen tekenreeksen. Zo betekent $!A \ident (!B \lor !C)$ dat de
tekenreeks~$!A$ gelijk is aan de tekenreeks die ontstaat door in deze
volgorde een openingshaakje, de tekenreeks $!B$, het symbool $\lor$, de
tekenreeks~$!C$ en een sluithaakje te concateneren. In dat geval weten
we dat het eerste symbool van~$!A$ een openingshaakje is, dat $!A$ de
tekenreeks $!B$ als deelreeks bevat (vanaf het tweede symbool), dat die
deelreeks wordt gevolgd door~$\lor$, enzovoort.

Omdat termen en !!{formula}s via inductieve definities uit basiselementen
worden opgebouwd, kunnen we met de volgende inductieprincipes
eigenschappen ervan bewijzen.

\begin{lem}[\emph{Inductieprincipe voor termen}]
    \ollabel{lem:trmind}
    Zij $\Lang L$ een eerste-ordetaal.
    Als een eigenschap~$P$ zo is dat
    %
    \begin{enumerate}
        \item zij geldt voor iedere !!{variable}~$v$,
        %
        \item zij geldt voor iedere !!{constant}~$a$ van~$\Lang L$, en
        %
        \item zij voor $f(t_1,\dotsc,t_n)$ geldt wanneer zij voor
        $t_1$,~\dots, $t_n$ geldt en $f$~een $n$-plaatsige
            !!{function} van~$\Lang L$ is,
    \end{enumerate}
    waarbij $t_1$,~\dots, $t_n$ termen van~$\Lang{L}$ zijn, dan geldt
    $P$ voor iedere term in~$\Trm[L]$.
\end{lem}

\begin{prob}
    Bewijs \olref[fol][syn][frm]{lem:trmind}.
\end{prob}

\begin{lem}[\emph{Inductieprincipe voor !!{formula}s}]
    \ollabel{thm:frmind}
    Zij $\Lang L$ een eerste-ordetaal.
    Als een eigenschap~$P$ voor alle atomaire !!{formula}s geldt
    en zo is dat
    %
    \begin{enumerate}
        \tagitem{prvNot}{zij voor $\lnot !A$ geldt wanneer zij
            voor~$!A$ geldt;}{}
        \tagitem{prvAnd}{zij voor $(!A \land !B)$ geldt
            wanneer zij voor $!A$ en~$!B$ geldt;}{}
        \tagitem{prvOr}{zij voor $(!A \lor !B)$ geldt
            wanneer zij voor $!A$ en~$!B$ geldt;}{}
        \tagitem{prvIf}{zij voor $(!A \lif !B)$ geldt
            wanneer zij voor $!A$ en~$!B$ geldt;}{}
        \tagitem{prvIff}{zij voor $(!A \liff !B)$ geldt
            wanneer zij voor $!A$ en~$!B$ geldt;}{}
        \tagitem{prvEx}{zij voor $\lexists[x][!A]$ geldt
            wanneer zij voor~$!A$ geldt;}{}
        \tagitem{prvAll}{zij voor $\lforall[x][!A]$ geldt
            wanneer zij voor~$!A$ geldt;}{}
  \end{enumerate}
  waarbij $!A$ en $!B$ !!{formula}s van~$\Lang{L}$ zijn, dan geldt
  $P$ voor alle formules in~$\Frm[L]$.
\end{lem}
\end{document}
